\section{A KKL theorem for biased Gaussian vectors: the proof of Theorem~\ref{t.KMSnonproduct}}\label{s.KMS}

In this section we prove Theorem~\ref{t.KMSnonproduct}. The proofs presented here do not rely on any results of the other sections. Recall that we use the following definition for influence. Given a vector $v\in\R^n$ and a Borel probability measure $\mu$ on $\R^n$ the influence of $v$ on $A$ under $\mu$ is
\[
I_{v,\mu}(A)=\underset{r \downarrow 0}{\liminf}\frac{\mu \left(A+[-r,r]v \right) - \mu \left(A \right)}{r} \in [0,+\infty]\,.
\]
\subsection{Sub-linearity of influences implies Theorem~\ref{t.KMSnonproduct}}\label{ss.KMSnonproduct}

The aim of this subsection is to prove Theorem~\ref{t.KMSnonproduct}. That this theorem holds for product Gaussian measures is a result by Keller, Mossel and Sen~\cite{keller2012geometric} (see Corollary~\ref{c.KMS} below). In order to extend this result to general Gaussian measures, we use a sub-linearity property for influences. In the present subsection, we state this property and use it to derive Theorem~\ref{t.KMSnonproduct} from the product case. In Subsection~\ref{ss.monotonic}, we prove the sub-linearity property for monotonic sets.

\paragraph{A KKL theorem for product Gaussian vectors} The authors of~\cite{keller2012geometric} introduce and study the notion of \emph{geometric influences} which are defined as follows: Let $\nu$ be a probability measure on $\R$ and let $A$ be a Borel subset of $\R^n$. If $i \in \lbrace 1,\dots, n \rbrace$, let:
\[
A_i^x := \lbrace y \in \R \, : \, (x_1,\dots, x_{i-1}, y, x_{i+1},\dots, x_n) \in A \rbrace \,.
\]
The geometric influence of $i$ on $A$ under the measure $\nu^{\otimes n}$ is:
\[
I^\mathcal{G}_{i,\nu}(A) := \E_{x \sim \nu^{\otimes n}} \left[\nu^+(A_i^x) \right] \in [0,+\infty] \,,
\]
where $\nu^+$ is the lower Minkowski content, defined as follows: for all $B \subseteq \R$ Borel,
\[
\nu^+(B) := \underset{r \downarrow 0}{\liminf} \frac{\nu\left(B+[-r,r]\right) - \nu \left(B\right)}{r} \in [0,+\infty] \,.
\]


In the case where $\mu=\nu^{\otimes n}$, $I^\calG_{i,\nu}$ and $I_{i,\mu}$ are closely related. Indeed, firstly, by Fubini's Theorem and Fatou's lemma, for each Borel subset $A\subset\R^n$ and each $i\in\{1,\dots,n\}$,
\begin{equation}\label{e.KMS.Fatou}
I_{i,\nu^{\otimes n}}(A)=\liminf_{r\downarrow 0} \E_{x \sim \nu^{\otimes n}} \left[\frac{\nu(A_i^x+[-r,r]) - \nu(A)}{r} \right]\geq I^\calG_{i,\nu}(A) \,.
\end{equation}

While the reverse inequality seems not true in general, we expect it to hold for a wide class of events. In particular, our Lemma~\ref{l.conv} and~(2.6) from~\cite{keller2012geometric} imply that this is the case for monotonic events. Since it is not useful to us, we do not investigate this matter any further.

We will need the following result, which is a direct consequence of Item~(1) of Theorem~1.5 of~\cite{keller2012geometric}. Several results of this type can also be found in the more recent~\cite{cordero2012hypercontractive} (see for instance the paragraph above Corollary~7 therein).

\begin{theo}\label{t.KMS}
There exists an absolute constant $c > 0$ such that the following holds:

Let $\nu = \mathcal{N}(0,1)$ and let $A$ be a Borel measurable subset of $\R^n$. Then:
\[
\sum_{i=1}^n I^\mathcal{G}_{i,\nu}(A) \geq c \, \nu^{\otimes n}(A) \, (1-\nu^{\otimes n}(A)) \, \sqrt{\log_+ \left(\frac{1}{\max_{i \in \lbrace 1,\dots, n \rbrace} I^\mathcal{G}_{i,\nu}(A)}\right) } \,.
\]
\end{theo}

\begin{coro}\label{c.KMS}
Let $\nu$ and $A$ be as in Theorem~\ref{t.KMS}. Then:
\[
\sum_{i=1}^n I_{i,\nu^{\otimes n}}(A) \geq c \, \nu^{\otimes n}(A) \, (1-\nu^{\otimes n}(A)) \, \sqrt{\log_+ \left(\frac{1}{\max_{i \in \lbrace 1,\dots, n \rbrace} I_{i,\nu^{\otimes n}}(A)}\right) } \,.
\]
\end{coro}

\begin{proof}
This is a direct consequence of Theorem~\ref{t.KMS} and Equation~\eqref{e.KMS.Fatou}.
\end{proof}

Now, let us state a sub-linearity property for the influences that we will be proved in Subsection~\ref{ss.monotonic}.

\begin{prop}\label{p.sublin}
Let $\Sigma$ be a $n \times n$ symmetric positive definite matrix, let $\mu \sim \mathcal{N}(0,\Sigma)$ and let $\left(e_1,\dots, e_n \right)$ be the canonical basis of $\R^n$. Also, let $v = \sum_{j=1}^n v_i \, e_i \in \R^n$. For every $i \in \lbrace 1,\dots, n \rbrace$ and every $A$ monotonic Borel subset of $\R^n$, we have:
\begin{equation}\label{e.sublinrec}
I_{v,\mu}(A) \leq \sum_{i=1}^n |v_i| \, I_{i,\mu}(A) \,.
\end{equation}
\end{prop}

We expect the above result to hold for a larger class of Borel subsets $A$ (in particular, we have not found any examples of Borel sets for which this does not hold) and not only for the decomposition on the canonical basis. See~\cite[Chapter~7, Appendix~A]{alelele} or~\cite[Chapter~2, Appendix~A]{hugogogo} for a proof of the inequality $I_{u+v,\mu}(A)\leq I_{u,\mu}(A)+I_{v,\mu}(A)$ for any $v,w$ when $A$ is a semi-algebraic set.

Let us now derive Theorem~\ref{t.KMSnonproduct} from Corollary~\ref{c.KMS} using Proposition\ref{p.sublin}.

\begin{proof}[Proof of Theorem~\ref{t.KMSnonproduct}] Let $\sqrt{\Sigma}$ be a symmetric square root of $\Sigma$ and let $\nu=\mathcal{N}(0,1)$. Then, $\mu=\calN(0,\Sigma)$ is the pushforward measure of $\nu^{\otimes n}$ by $\sqrt{\Sigma}$. Thanks to Corollary~\ref{c.KMS} (applied to the event $\sqrt{\Sigma}^{-1}(A)$), it is sufficient to prove the following claim.
\end{proof}

\begin{enonce}{Claim}
Let $\nu = \mathcal{N}(0,1)$. We have:
\begin{equation}\label{e.cons_sublin_1}
\max_{j\in\lbrace 1,\dots, n\rbrace} I_{j,\nu^{\otimes n}}(\sqrt{\Sigma}^{-1} \left(A \right)) \leq \| \sqrt{\Sigma} \|_{\infty,op} \cdot \max_{i \in \lbrace 1,\dots, n \rbrace} I_{i,\mu}(A) \,.
\end{equation}
Moreover:
\begin{equation}\label{e.cons_sublin_2}
\sum_{j=1}^n I_{j,\nu^{\otimes n}}(\sqrt{\Sigma}^{-1} \left(A \right)) \leq \| \sqrt{\Sigma} \|_{\infty,op} \cdot \sum_{i=1}^n I_{i,\mu}(A) \,.
\end{equation}
\end{enonce}

\begin{proof}
For each $j\in\lbrace 1,\dots, n\rbrace$:
\begin{align*}
I_{j,\nu^{\otimes n}}(\sqrt{\Sigma}^{-1} \left(A \right)) & = \underset{r \downarrow 0}{\liminf} \frac{\nu^{\otimes n} \left(\sqrt{\Sigma}^{-1} \left(A \right) + [-r,r] \, e_j \right) - \nu^{\otimes n} \left(\sqrt{\Sigma}^{-1} \left(A \right) \right)}{r}\\
& = \underset{r \downarrow 0}{\liminf} \frac{\mu \left(A + [-r,r] \, \sqrt{\Sigma} \cdot e_j \right) - \mu \left(A \right)}{r}\\
& = I_{\sqrt{\Sigma} \cdot e_j,\mu}(A) \,.
\end{align*}
By using Proposition~\ref{p.sublin}, we obtain that:
\begin{equation}\label{e.cons_sublin_inter}
I_{j,\nu^{\otimes n}}(\sqrt{\Sigma}^{-1} \left(A \right)) \leq \sum_{i = 1}^n |\sqrt{\Sigma}(i,j)| I_{i,\mu}(A) \,.
\end{equation}
Hence (since $\sqrt{\Sigma}$ is symmetric):
\begin{align*}
I_{j,\nu^{\otimes n}}(\sqrt{\Sigma}^{-1} \left(A \right))& \leq \sum_{i = 1}^n |\sqrt{\Sigma}(i,j)| \max{i \in \lbrace 1,\dots, n \rbrace} I_{i,\mu}(A) \\
&\leq \|\sqrt{\Sigma} \|_{\infty,op} \, \max_{i \in \lbrace 1,\dots, n \rbrace} I_{i,\mu}(A) \,.
\end{align*}
We obtain~\eqref{e.cons_sublin_1} by taking the supremum over $j$. Inequality~\eqref{e.cons_sublin_inter} also implies that:
\begin{align*}
\sum_{j=1}^n I_{\sqrt{\Sigma} \cdot e_j,\mu}(A) & \leq \sum_{j=1}^n \sum_{i=1}^n |\sqrt{\Sigma}(i,j)| \, I_{i,\mu}(A)
 = \sum_{i=1}^n I_{i,\mu}(A) \sum_{j=1}^n |\sqrt{\Sigma}(i,j)|\\
& \leq \|\sqrt{\Sigma}\|_{\infty,op} \cdot \sum_{i=1}^n I_{i,\mu}(A) \,,
\end{align*}
which is~\eqref{e.cons_sublin_2}.
\end{proof}

\subsection{Sub-linearity of influences for monotonic events}\label{ss.monotonic}

The proof of Proposition~\ref{p.sublin} relies on the following lemma.

\begin{lemm}\label{l.conv}
Let $\Sigma$ be a $n \times n$ symmetric positive definite matrix and let $\mu = \mathcal{N}(0,\Sigma)$. Moreover, let $v \in \R^n$. For every $A$ monotonic Borel subset of $\R^n$ and every $i \in \lbrace 1,\dots, n \rbrace$, we have:
\[
\exists \lim_{r \downarrow 0} \frac{\mu \left(A + [-r,r] e_i + [-r,r]v \right) - \mu \left(A + [-r,r]v \right)}{r} = I_{i,\mu}(A) \,.
\]
\end{lemm}

We first prove Proposition~\ref{p.sublin} using Lemma~\ref{l.conv}.

\begin{proof}[Proof of Proposition~\ref{p.sublin}]
Let $A$ be a monotonic Borel subset of $\R^n$ and fix $v \in \R^n$. Let $v^i=\sum_{k=i}^n v_k e_k$. We will prove that, for every $i \in \{ 1,\dots, n-1 \}$:
\begin{equation}\label{e.sublin}
I_{v^i,\mu}(A) \leq |v_i| \, I_{i,\mu}(A) + I_{v^{i+1},\mu}(A) \,.
\end{equation}
The result will then follow directly by induction (and since $I_{v^{n},\mu}(A)= I_{n,\mu}(A)$). For any $w_1,w_2 \subseteq \R^n$, we have $[-r,r](w_1+w_2) \subseteq [-r,r]w_1+[-r,r]w_2$. Hence:
\begin{align*}
\mu \left(A + [-r,r]v^i \right) &=\mu \left(A + [-r,r] \left(v_i e_i + v^{i+1} \right) \right)\\
&\leq \mu \left(A + [-|v_i| r,|v_i| r] e_i + [-r,r] v^{i+1} \right)\\
&= \mu \left(A + [-|v_i| r,|v_i| r] e_i + [-r,r] v^{i+1} \right)\\
&\qquad - \mu \left(A + [-r,r] v^{i+1} \right) + \mu \left(A + [-r,r] v^{i+1} \right) \,.
\end{align*}

By Lemma~\ref{l.conv} we have:
\[
\frac{\mu \left(A + [-|v_i| r,|v_i| r] \, e_i + [-r,r] \, v^{i+1} \right) - \mu \left(A + [-r,r]v^{i+1} \right)}{r} \underset{r \downarrow 0}{\longrightarrow} |v_i| \, I_{i,\mu}(A) \,.
\]
Equation~\eqref{e.sublin} follows.
\end{proof}

\begin{proof}[Proof of Lemma~\ref{l.conv}] We are inspired by the proof of Proposition~1.3 in~\cite{keller2012geometric}. We write the proof for $A$ decreasing since the proof for $A$ increasing is identical. Also, we prove the result in the case $i = n$. We write $\tilde{x}$ for the first $(n-1)$ coordinates of any $x\in\R^n$. Let $A_r = A + [-r,r]v$, note that $A_r$ is decreasing, and write for any $\tilde{x}\in\R^{n-1}$:
\begin{align*}
s(\tilde{x}) &:= \sup \lbrace x_n \in \R \, : \, (\tilde{x}, x_n) \in A \rbrace \in [-\infty,+\infty] \,,\\
s_r(\tilde{x}) &:= \sup \lbrace x_n \in \R \, : \, (\tilde{x}, x_n) \in A_r \rbrace \in [-\infty,+\infty] \,.
\end{align*}
Let $\lambda$ be the density function of $\mu$. Since $A$ and $A_r$ are decreasing, we have:
\begin{align}\label{e.mon.integral}
\frac{\mu (A + [-r,r] e_n) - \mu (A)}{r} &= \frac{1}{r} \int_{\R^{n-1}} \left(\int_{s(\tilde{x})}^{s(\tilde{x})+r} \lambda(\tilde{x},x_n) \dd x_n\right)\dd\tilde{x} \,, \nonumber\\
\frac{\mu (A_r + [-r,r] e_n) - \mu (A_r)}{r} &= \frac{1}{r} \int_{\R^{n-1}} \left(\int_{s_r(\tilde{x})}^{s_r(\tilde{x})+r} \lambda(\tilde{x},x_n) \dd x_n\right)\dd\tilde{x} \,,
\end{align}
where by convention $\int_{-\infty}^{-\infty+r} = \int_{+\infty}^{+\infty+r} = 0$. For each $\tilde{x}\in\R^{n-1}$ let:
\[
g_1(\tilde{x})=\sup_{x_n\in\R}\lambda(\tilde{x},x_n);\: \: g_2(\tilde{x})=\sup_{x_n\in\R}\left|\frac{\partial\lambda}{\partial x_n}(\tilde{x},x_n)\right|.
\]
Direct computation shows that $g_1,g_2\in L^1(\R^{n-1})$. By the mean value inequality, for each $\tilde{x}\in\R^{n-1}$:
\[
\left|\frac{1}{r}\left(\int_{s(\tilde{x})}^{s(\tilde{x})+r} \lambda(\tilde{x},x_n) \dx_n\right)-\lambda(\tilde{x},s(\tilde{x}))\right|+\left|\frac{1}{r}\left(\int_{s_r(\tilde{x})}^{s_r(\tilde{x})+r} \lambda(\tilde{x},x_n) \dd x_n\right)-\lambda(\tilde{x},s_r(\tilde{x}))\right|
\]
is no greater than $2rg_2(\tilde{x})$. Combining this with Equation~\eqref{e.mon.integral} we get:
\begin{multline*}
\left|\frac{\mu (A + [-r,r] e_n) - \mu (A)}{r}- \frac{\mu (A_r + [-r,r] e_n) - \mu (A_r)}{r} \right|\\
\leq \left|\int_{\R^{n-1}}\lambda(\tilde{x},s(\tilde{x})) -\lambda(\tilde{x},s_r(\tilde{x})) \, \dd\tilde{x}\right|+2r\int_{\R^{n-1}}g_2(\tilde{x}) \, \dd\tilde{x}\,.
\end{multline*}
Since $g_2\in L^1(\R^{n-1})$, the second integral in the last inequality is finite and independent of $r$. Moreover:
\[
\int_{\R^{n-1}}\lambda(\tilde{x},s_r(\tilde{x})) -\lambda(\tilde{x},s(\tilde{x})) \, \dd\tilde{x}=\int_{\R^{n-1}}\int_{s(\tilde{x})}^{s_r(\tilde{x})}\frac{\partial\lambda}{\partial x_n}(\tilde{x},x_n) \, \dd x_n \dd\tilde{x}\,.
\]
Since $\frac{\partial\lambda}{\partial x_n}\in L^1(\R^n)$, by dominated convergence, all that remains is to show that for a.e. $\tilde{x}\in\R^{n-1}$: $s_r(\tilde{x}) \longrightarrow_{r\downarrow 0} s(\tilde{x})$. Since for each $\tilde{x}$ the sequence $s_r(\tilde{x})$ is decreasing, it converges to some $s_\infty(\tilde{x}) \geq s(\tilde{x})$. Let us prove that, for a.e. $\tilde{x}\in\R^{n-1}$, $s_\infty(\tilde{x})=s(\tilde{x})$. To do so, first note that, since $A$ is decreasing, we have:
\[
0 \leq \mu(A_r) - \mu(A) \leq \Pro \left[X - \sum_{i=1}^n r |v_i|e_i \in A \right] - \Pro \left[X \in A \right] \,,
\]
where $X \sim \calN(0,\Sigma)$. By dominated convergence, the right hand side tends to $0$ when $r\rightarrow 0$. Now, note that:
\[
\mu(A_r)-\mu(A) = \int_{\R^{n-1}} \int_{s(\tilde{x})}^{s_r(\tilde{x})} \lambda(\tilde{x},x_n) \, \dd x_n \dd\tilde{x} \,.
\]
Hence, by Fatou's lemma:
\[
0 = \int_{\R^{n-1}} \int_{s(\tilde{x})}^{s_\infty(\tilde{x})} \lambda(\tilde{x},x_n) \, \dd x_n \dd\tilde{x} \,.
\]
Since the $\lambda$ takes only positive values, this implies that for a.e. $\tilde{x}\in\R^{n-1}$, $s_\infty(\tilde{x})=s(\tilde{x})$.
\end{proof}
